
On 18\% of factuality questions, token entropy and answer consistency point in opposite directions---and each method is right about its own subset. This observation motivates a systematic investigation of whether existing uncertainty quantification approaches capture fundamentally different hallucination patterns.

Large language models generate fluent text that may be factually incorrect, a phenomenon termed hallucination. Detecting such errors before deployment is relevant for applications in domains where factual accuracy matters. Two approaches have emerged: entropy-based methods that measure uncertainty in the model's token probability distribution, and consistency-based methods that measure agreement across multiple sampled responses.

These approaches have been developed and evaluated independently. Kuhn et al. (2023) introduced semantic entropy for natural language generation tasks, while Manakul et al. (2023) developed SelfCheckGPT for consistency-based detection on biography generation. Neither work directly compared these methods on the same factuality benchmark with the same model. This leaves a gap: it is unknown whether entropy and consistency capture redundant information or complementary signals.

This work presents the first systematic comparison of token entropy and N-sample consistency on TruthfulQA with LLaMA-2-7B. The key finding is that these signals are orthogonal (r = 0.228), sharing only approximately 5\% of variance. The contributions are:

\begin{enumerate}
\item \textbf{Systematic comparison} of entropy and consistency methods on the same factuality benchmark (TruthfulQA) with the same model (LLaMA-2-7B).
\end{enumerate}

\begin{enumerate}
\item \textbf{Empirical demonstration of orthogonality} (r = 0.228, p < $10^{-10})$, showing these methods capture different aspects of model uncertainty.
\end{enumerate}

\begin{enumerate}
\item \textbf{Complementarity analysis} showing that 18.1\% of questions exhibit method disagreement, with each method achieving AUROC > 0.75 on its winning subset.
\end{enumerate}

